\chapter{Conclusion and Future Work}
\label{ch:conclusion}

This thesis has investigated the problem of formal-conformance fault prevention in
LLM-assisted Agile-SOFL development.  We designed and evaluated the HSP-Agile pipeline, a
four-stage hybrid system that combines LLM candidate synthesis, Z3-based formal
conformance checking, counterexample-guided repair, and pattern-guard screening under a
conjunctive acceptance predicate.  This chapter summarises the contributions and key
empirical findings, and charts directions for future research.

% ---------------------------------------------------------------------------
\section{Summary of Contributions}
\label{sec:conc:contributions}
% ---------------------------------------------------------------------------

This thesis makes five primary contributions to the literature on LLM-assisted formal
software development within the Agile-SOFL framework~\citep{li2022agilesofl,liu2014sofl}.

\paragraph{Contribution 1: The HSP-Agile Pipeline.}
We designed, implemented, and evaluated the HSP-Agile pipeline end-to-end.  The pipeline
comprises four stages: (1)~LLM Candidate Synthesis, which queries a large language model
with an FSF specification and optional repair feedback; (2)~Formal Conformance Checker,
which employs Z3-guided test-case generation~\citep{demoura2008z3,barrett2018smt} and oracle comparison to compute a strict
conformance score $\sigma(c, S)$; (3)~Counterexample-Guided Repair, which constructs
targeted repair prompts from witnessed test failures~\citep{solarlezama2008sketching,lynn2024verifierloop}; and (4)~Pattern Guard, which screens
candidates against a catalogue of high-severity Agile-SOFL fault patterns~\citep{li2023iet,hovemeyer2004findbugs}.  The pipeline
is parameterised by a conformance threshold $\tau$ and a maximum repair budget $K$, with
a graceful fallback to the best-scoring candidate when the budget is exhausted.  The
complete implementation, benchmark, and experiment artefacts are made publicly available
to support replication.

\paragraph{Contribution 2: The Strict Acceptance Predicate.}
We formally defined the strict acceptance predicate $\mathcal{A}(c, S)$ as the conjunction
of a conformance threshold gate ($\sigma(c, S) \geq \tau$) and a pattern guard gate
($\neg \exists p \in \mathcal{P} : \text{match}(c, p)$).  This predicate provides a
principled and operationalisable acceptance criterion that goes beyond the binary pass/fail
criteria used by existing LLM code-generation benchmarks.  We showed empirically that the
binary all-or-nothing strict success rate is an insufficient discriminator on hard
precedence-sensitive benchmarks, motivating the use of formal conformance as a richer
primary metric.

\paragraph{Contribution 3: The Hard Benchmark.}
We constructed a benchmark of 120 Agile-SOFL tasks, all drawn from the hard partition of
precedence-sensitive FSF specifications.  The benchmark is designed to maximally stress
scenario-ordering conformance: all tasks require the synthesised candidate to correctly
implement at least one non-trivial precondition dependency between scenario branches.
B0 scores of 27.5\% confirm that the benchmark is genuinely hard, not solvable by naive
template generation.  The benchmark is released as a standalone artefact for use by future
researchers evaluating LLM-assisted formal synthesis tools~\citep{loughman2024dafnybench}.

\paragraph{Contribution 4: Two-Layer Prevention Protocol.}
We proposed and partially evaluated a two-layer mutation-based prevention protocol
measuring probabilistic detection rate (PDR) and false accept rate (FAR).  The protocol
injects seeded conformance faults into known-correct candidates and measures whether each
pipeline mode correctly rejects the mutant.  Preliminary results across 40 tasks per mode
(B1: PDR 7.5\%, B2: PDR 15.0\%, M: PDR 10.0\%) establish a baseline for future full-scale
prevention analysis.  The protocol itself is a methodological contribution: it provides a
rigorous framework for evaluating fault-prevention pipelines that is independent of the
specific pipeline design.

\paragraph{Contribution 5: Reproducible Evaluation Methodology.}
The evaluation spans 840 experiment runs (7 modes $\times$ 120 tasks) with fixed random
seeds, deterministic Z3 solving, and a standardised result schema.  All runs are logged
with full audit traces including per-task conformance scores, per-repair-iteration
conformance trajectories, pattern guard outputs, and raw LLM responses.  Statistical
analysis uses Wilcoxon signed-rank tests with two-tailed $p$-values~\citep{ammann2016testing}.  This reproducibility
infrastructure is released alongside the pipeline implementation to enable independent
replication and extension.

% ---------------------------------------------------------------------------
\section{Key Empirical Findings}
\label{sec:conc:findings}
% ---------------------------------------------------------------------------

The experimental evaluation yields five findings of central importance.

\paragraph{Finding 1: HSP-Agile improves formal conformance by 6.2 percentage points over
the one-shot baseline.}
The full pipeline M achieves a mean strict formal conformance of \textbf{90.4\%} on the
120-task hard benchmark.  This represents a gain of \textbf{+6.2 percentage points} over
the one-shot LLM baseline B1 (84.2\%) and \textbf{+2.4 percentage points} over the
test-feedback baseline B2 (88.0\%).  These gains are consistent across the full task set
and are not driven by a small number of easy tasks.

\paragraph{Finding 2: The repair loop is the dominant contributor to conformance
improvement.}
Ablation A3, which removes the counterexample-guided repair loop while retaining the
formal checker and pattern guard, reduces mean conformance to \textbf{84.1\%}, a drop of
\textbf{$-$6.3 percentage points} relative to M.  This makes the repair loop the single
most important component in the pipeline: its removal costs more than returning to the
one-shot baseline.  The result validates the design decision to invest in repair-prompt
construction from Z3 witnesses rather than simply re-querying the LLM without targeted
feedback.

\paragraph{Finding 3: The formal conformance checker contributes independently.}
Ablation A1, which removes the formal checker while retaining the repair loop (which then
receives only unit-test feedback) and the pattern guard, reduces mean conformance to
\textbf{86.5\%}, a drop of \textbf{$-$3.9 percentage points} relative to M.  This confirms
that the Z3-guided test cases provide qualitatively different and complementary information
to the unit-test feedback used by B2.

\paragraph{Finding 4: Latency overhead is the primary cost of the pipeline.}
Mode M incurs a mean latency of \textbf{43{,}659\,ms} per task, compared with
\textbf{10{,}769\,ms} for B2 �?a \textbf{4.6$\times$ overhead}.  This overhead is
attributable to the increased number of LLM API calls (mean 2.08 per task under M) and
the Z3 conformance-checking steps.  The overhead is the principal trade-off introduced by
HSP-Agile relative to simpler baselines, and its practical acceptability depends on the
deployment context (see Section~\ref{sec:disc:implications}).

\paragraph{Finding 5: Strict success rate is insufficient to discriminate pipeline
quality on hard benchmarks.}
A Wilcoxon signed-rank test comparing M and B2 on strict success rate yields
$p = 0.921$, which is not statistically significant.  M's strict success rate of 25.0\%
is marginally \emph{lower} than B2's 26.7\%, yet M achieves higher formal conformance.
This paradox, analysed in detail in Section~\ref{sec:disc:paradox}, arises because M's
conjunctive acceptance gate rejects some candidates that pass all test cases but
trigger pattern guard alerts.  Binary pass/fail metrics fail to capture the quality
improvement that HSP-Agile delivers; formal conformance is the appropriate primary metric
for evaluating hard-benchmark performance.

% ---------------------------------------------------------------------------
\section{Future Research Directions}
\label{sec:conc:future}
% ---------------------------------------------------------------------------

The findings of this thesis open several avenues for future investigation.

\paragraph{Richer specification support.}
The current HSP-Agile pipeline targets FSF specifications with scenario-ordering
constraints expressible as Z3 integer constraints.  Many real Agile-SOFL specifications
additionally involve temporal properties (e.g.\ ``scenario A must always precede scenario
B within the same session''), stateful behaviours (e.g.\ modifications to shared database
tables), and composite output structures.  Extending the formal checker to handle these
richer specification classes �?for example, by incorporating linear temporal logic (LTL)
model checking or relational database constraint solving �?would substantially broaden the
scope of FSF tasks that HSP-Agile can support.  The challenge is to do so without
disproportionately increasing latency.

\paragraph{Theorem-prover backends for stronger guarantees.}
The Z3-based checker provides bounded guarantees relative to a finite generated test suite.
A natural extension is to replace or augment the Z3 checker with a theorem-prover backend
capable of generating complete correctness certificates.  Dafny~\citep{leino2010dafny} and related
deductive-verification ecosystems~\citep{loughman2024dafnybench,woodcock2009formalsurvey} are promising candidates: they provide stronger
guarantees on accepted candidates while retaining the repair loop's ability to guide the
LLM toward correct implementations.

\paragraph{Adaptive budget allocation.}
The current pipeline uses a fixed repair budget $K$ for all tasks.  This is suboptimal:
easy tasks may exhaust the budget on unnecessary iterations, while hard tasks may be
abandoned just before convergence.  An adaptive budget-allocation strategy �?for example,
one that halts early when conformance improvement per iteration falls below a threshold,
and allocates saved budget to tasks that show consistent improvement �?could reduce mean
latency while maintaining or improving conformance on the hardest tasks.  Reinforcement
learning over iteration trajectories is a natural framing for this optimisation.

\paragraph{Industrial replication.}
The benchmark used in this thesis is synthetic.  Industrial replication is essential before
the quantitative findings can be generalised.  We propose collaborating with Agile-SOFL
practitioner groups to obtain real FSF specifications from software development projects,
applying HSP-Agile to LLM-assisted implementation of those specifications, and measuring
conformance against the formally specified acceptance criteria used in those projects.
Industrial tasks will likely differ from the benchmark in specification complexity, scenario
count, and the prevalence of the specific fault patterns catalogued in the pattern guard.

\paragraph{Model comparison and prompt engineering.}
All experiments in this thesis used Qwen2.5-27B-Instruct via an OpenAI-compatible LLM gateway.  A systematic comparison across
state-of-the-art LLMs---GPT-4o, Claude~3.5 Sonnet, Gemini~1.5 Pro---on the same 120-task
hard benchmark would establish whether the conformance gains delivered by HSP-Agile are
model-specific or model-agnostic~\citep{openai2023gpt4,chen2021codex,dakhel2023copilot}.  Such a comparison would also enable a systematic study
of how repair prompt design interacts with model-specific instruction-following
capabilities.  We anticipate that models with stronger instruction-following would exhibit
faster repair convergence, reducing both latency and budget consumption.

\paragraph{Full prevention evaluation.}
The mutation-based PDR/FAR evaluation reported in this thesis is preliminary (40 tasks per
mode).  A full evaluation should span all 120 tasks, all seven modes, and a richer mutant
population including multi-fault mutants (candidates with two or more injected ordering
violations).  A full evaluation would also include false-accept rate analysis across all
modes, enabling direct comparison of the prevention quality of M against B1, B2, and the
ablations.

\bigskip
\noindent
In summary, HSP-Agile demonstrates that coupling LLM synthesis with lightweight formal
conformance checking and counterexample-guided repair produces measurable and reproducible
quality improvements in Agile-SOFL code generation.  The pipeline is not a replacement for
full formal verification, but it provides a practical bridge between informal agile
development practices and the formal quality standards demanded by downstream verification
artefacts �?a bridge built from bounded formal checking, targeted repair, and principled
conjunctive acceptance.
